% ------------------------------------------------------------------
% Chapter 3 — Requirements (Jarvis)
% Converted from docs/requirements.md for the UCL MSc CS project report.
% Compile standalone with pdflatex/xelatex/tectonic, or \input into the
% main report (strip the preamble and \setcounter when doing so).
% ------------------------------------------------------------------
\documentclass[12pt]{article}

\usepackage[a4paper,margin=2.5cm]{geometry}
\usepackage[T1]{fontenc}
\usepackage[utf8]{inputenc}
% Times New Roman text and matching maths; fall back to mathptmx if
% the newtx bundle is not available on this installation.
\IfFileExists{newtxtext.sty}{%
  \usepackage{newtxtext,newtxmath}%
}{%
  \usepackage{mathptmx}%
}
\usepackage{booktabs}
\usepackage{longtable}
\usepackage{tabularx}
\usepackage{array}
\usepackage{enumitem}
\usepackage{microtype}

% Requirement-text column that wraps and stays ragged-right.
\newcolumntype{R}[1]{>{\raggedright\arraybackslash}p{#1}}

% This file is Chapter 3 of the report: make \section come out as 3.
\setcounter{section}{2}

\begin{document}

\section{Requirements}\label{sec:requirements}

This chapter follows the department's convention: a requirement-gathering
narrative, personas, MoSCoW-prioritised functional and non-functional
requirement tables, use cases, and user flows. Requirements state
\textbf{what the user can do and what the system must guarantee} --- how any
of it is achieved belongs to the Design chapter, not here.

\subsection{Problem statement}\label{sec:problem-statement}

Capable AI assistants exist, but the ones available to consumers are
cloud-hosted: whatever the user asks --- their email, their calendar, their
files --- leaves their machine. Assistants that do run locally exist as
developer tools: they require command-line installation, manual model
management, and constant configuration, placing them out of reach of the
people who would benefit most. And in either form, today's assistants have a
fixed repertoire: a request outside it simply fails, or worse, is attempted
badly and reported as done.

There is a need for a personal assistant that a non-technical person can
install and operate on their own computer, that keeps every piece of their
data on that computer, that handles the everyday load of email and calendar,
that can safely \emph{extend its own abilities} when asked for something new
--- and that never takes an irreversible action, or installs code into
itself, without the user's knowledge and consent.

\subsection{Requirement gathering}\label{sec:requirement-gathering}

Direct stakeholder interviews were not feasible; requirements were derived
from four sources:

\begin{itemize}[itemsep=2pt]
  \item \textbf{The project brief.} The supervisors' proposal specifies the
    pillars: a locally-run, privacy-first assistant; autonomous skill
    generation in response to novel requests; asynchronous access through a
    familiar messaging channel; and evaluation of latency, task success, and
    the security of sandboxed execution.
  \item \textbf{Measured behaviour of an existing agent runtime.} A pilot
    study run for this project (Chapter~5) drove a mature open-source agent
    framework with local models on consumer hardware. Its failure modes ---
    actions reported done that never happened, uncontrolled tool use,
    self-written capabilities installed without testing --- directly motivate
    the honesty, consent, and verification requirements below.
  \item \textbf{Community evidence.} Published community guidance on running
    agent frameworks against local models corroborates the pilot: reliable
    operation is generally reported only with far larger models than consumer
    hardware can serve, reinforcing the need for a system designed around
    small-model limits rather than assuming them away.
  \item \textbf{Scenario analysis.} Personas and scenarios
    (\S\ref{sec:personas}) were mapped against the everyday tasks a personal
    assistant is expected to carry: correspondence, scheduling, file chores,
    and repeated ad-hoc requests.
\end{itemize}

\subsubsection{Personas}\label{sec:personas}

\begin{table}[htbp]
  \centering
  \caption{Persona: Maya Okafor}
  \label{tab:persona-maya}
  \begin{tabularx}{\textwidth}{@{}>{\bfseries}p{3.4cm}X@{}}
    \toprule
    \multicolumn{2}{@{}l}{\textbf{Maya Okafor}} \\
    \midrule
    Role & Self-employed accountant, 52 \\
    Background & Runs her practice from a MacBook; competent with ordinary
      applications, has never used a terminal. Client correspondence is
      confidential by professional obligation. \\
    Responsibilities & Managing a high volume of client email, deadlines and
      appointments; safeguarding client financial data. \\
    Challenges & Maya would benefit most from an assistant, but every capable
      one she has tried requires sending client data to a third-party
      service, which she cannot justify professionally. Local alternatives
      assume technical skill she does not have and does not want to acquire.
      She needs installation and daily use to feel like any other Mac
      application, and she needs certainty that nothing she asks ever leaves
      her machine. \\
    \bottomrule
  \end{tabularx}
\end{table}

\begin{table}[htbp]
  \centering
  \caption{Persona: Dev Sharma}
  \label{tab:persona-dev}
  \begin{tabularx}{\textwidth}{@{}>{\bfseries}p{3.4cm}X@{}}
    \toprule
    \multicolumn{2}{@{}l}{\textbf{Dev Sharma}} \\
    \midrule
    Role & MSc student, 24 \\
    Background & Technically literate; comfortable trying new software;
      time-poor. \\
    Responsibilities & Coursework deadlines, supervisor correspondence,
      part-time work scheduling. \\
    Challenges & Dev's questions are small but constant --- ``did the
      department reply about my extension?'', ``what's on Friday?'' --- and
      each one costs a context switch into webmail. He wants them answered
      conversationally, ideally by voice while doing something else, and
      wants repeated chores to get faster rather than cost the same effort
      every time. \\
    \bottomrule
  \end{tabularx}
\end{table}

\begin{table}[htbp]
  \centering
  \caption{Persona: Ines Almeida}
  \label{tab:persona-ines}
  \begin{tabularx}{\textwidth}{@{}>{\bfseries}p{3.4cm}X@{}}
    \toprule
    \multicolumn{2}{@{}l}{\textbf{Ines Almeida}} \\
    \midrule
    Role & Software developer, 34 \\
    Background & Builds backend systems professionally; security-conscious by
      habit and by trade. \\
    Responsibilities & Evaluating tools before trusting them with real data. \\
    Challenges & Ines's default stance toward an agent that can act on her
      machine is distrust. She will not use an assistant that acts invisibly:
      she wants to see what it is doing while it does it, review what any
      self-installed ability is permitted to touch, veto anything
      irreversible, and verify --- not be told --- that claimed outcomes
      really happened. \\
    \bottomrule
  \end{tabularx}
\end{table}

\paragraph{Scenario.} Maya asks the assistant whether a client has replied
about a filing deadline; it answers from her own mailbox, on her own machine,
in seconds. She then asks it to rename three hundred scanned receipts by date
--- something it has never done before. It tells her it doesn't yet have that
ability, describes what it would build and what that ability would be allowed
to access, and asks permission. She approves; a minute later the receipts are
renamed, and the ability remains for next quarter. Ines, watching the same
flow, can open the activity view mid-task, read the new ability's
permissions, and delete it afterwards if she chooses.

\subsubsection{MoSCoW table}\label{sec:moscow}

\begin{longtable}{@{}l R{11.3cm} l@{}}
  \caption{Functional requirements (MoSCoW-prioritised)}
  \label{tab:functional-reqs} \\
  \toprule
  \textbf{ID} & \textbf{Functional Requirement} & \textbf{Priority} \\
  \midrule
  \endfirsthead
  \caption[]{Functional requirements (MoSCoW-prioritised) --- continued} \\
  \toprule
  \textbf{ID} & \textbf{Functional Requirement} & \textbf{Priority} \\
  \midrule
  \endhead
  \midrule
  \multicolumn{3}{r@{}}{\small\emph{Continued on next page}} \\
  \endfoot
  \bottomrule
  \endlastfoot
  F1 & The application shall allow the user to install and set it up without
    any command-line interaction & MUST \\
  F2 & The application shall allow the user to make requests in natural
    language through a chat interface & MUST \\
  F3 & The application shall allow the user to speak requests aloud and hear
    spoken replies & MUST \\
  F4 & The application shall allow the user to read, search and summarise
    email from their own account & MUST \\
  F5 & The application shall allow the user to draft, reply to and send email
    from their own account & MUST \\
  F6 & The application shall allow the user to ask what is on their calendar
    & MUST \\
  F7 & The application shall not take an irreversible action unless the
    user's request authorised it or the user explicitly approves it when
    asked & MUST \\
  F8 & The application shall offer to build a new skill when a request is
    beyond its abilities, and shall proceed only with the user's explicit
    approval & MUST \\
  F9 & The application shall test every skill it builds and shall install
    only those that pass & MUST \\
  F10 & The application shall show the user every ability it has, including
    what each is permitted to access & MUST \\
  F11 & The application shall allow the user to remove any learned ability
    & MUST \\
  F12 & The application shall show the user what it is doing while it works
    & MUST \\
  F13 & The application shall report outcomes from verified results and shall
    report failures as failures & MUST \\
  F14 & The application shall not enter, store or ask for the user's
    passwords or payment details, and shall hand control to the user at any
    sign-in & MUST \\
  F15 & The application shall allow the user to choose which
    locally-installed models it uses & SHOULD \\
  F16 & The application shall make skills it has built usable by other agent
    software on the user's machine & SHOULD \\
  F17 & The application shall allow the user to view its learned web routines
    and remove them & SHOULD \\
  F18 & The application shall allow the user to review past skill-building
    attempts, including failed ones & SHOULD \\
  F19 & The application shall allow the user to view measurements of its own
    performance & SHOULD \\
  F20 & The application shall allow the user to have appointments and events
    booked into their calendar --- either as requested directly, or from
    details the assistant finds in their email --- subject to the same
    approval rules as any other irreversible action (F7) & SHOULD \\
F28 & The application shall allow a request beyond its own abilities to be handed to other agent software on the user's machine, only with the user's explicit approval and with the handoff's reduced guarantees stated & SHOULD \\
  F21 & The application shall allow the user to interact with it from their
    phone & COULD \\
  F22 & The application shall allow the user to make quick requests without
    opening the main window & COULD \\
  F23 & The application shall allow the user to browse their past
    conversations & COULD \\
  F24 & The application shall allow the user to share a skill it has built
    with another installation & COULD \\
  F25 & The application shall automate websites beyond the user's mail and
    calendar & WON'T \\
  F26 & The application shall run on platforms other than Apple-silicon macOS
    & WON'T \\
  F27 & The application shall support multiple users on one installation
    & WON'T \\
\end{longtable}

\begin{longtable}{@{}l R{11.3cm} l@{}}
  \caption{Non-functional requirements (MoSCoW-prioritised)}
  \label{tab:nonfunctional-reqs} \\
  \toprule
  \textbf{ID} & \textbf{Non-Functional Requirement} & \textbf{Priority} \\
  \midrule
  \endfirsthead
  \caption[]{Non-functional requirements (MoSCoW-prioritised) --- continued} \\
  \toprule
  \textbf{ID} & \textbf{Non-Functional Requirement} & \textbf{Priority} \\
  \midrule
  \endhead
  \midrule
  \multicolumn{3}{r@{}}{\small\emph{Continued on next page}} \\
  \endfoot
  \bottomrule
  \endlastfoot
  NF1 & The application shall keep the user's data, and all processing of it,
    on the user's machine, and shall not transmit it to any third-party
    service & MUST \\
  NF2 & The application shall not allow content it encounters while working
    --- a webpage, an email, a document --- to authorise actions & MUST \\
  NF3 & The application shall be usable by a first-time, non-technical user
    without documentation & MUST \\
  NF4 & The application shall operate fully on a consumer machine with 24~GB
    of unified memory & MUST \\
  NF5 & The application shall acknowledge every request within a few seconds,
    answer routine questions in under a minute, and complete typical mail
    tasks within about two & MUST \\
  NF6 & The application shall recover from interruption gracefully, and shall
    never present a failed or unfinished action as a completed one & SHOULD \\
  NF7 & The application shall become faster at tasks it has performed before
    & SHOULD \\
  NF8 & The application shall be accompanied by documented code, a system
    manual and a user manual & SHOULD \\
\end{longtable}

\paragraph{Scope remarks.} F25 is excluded on evidence rather than
preference: the pilot measured the baseline runtime at 1/7 on general web
tasks with local models, and community guidance places reliable agentic
browsing at model sizes beyond consumer hardware; mail and calendar (F4--F6)
remain in scope because the pilot verified them end-to-end. F21 is held at
COULD as an inexpensive extension of the same request pipeline over a
messaging transport. The prohibitions are deliberately MUSTs rather than
WON'Ts: F7, F14 and NF1 are binding properties of every release, not features
deferred from this one.

\subsection{Use cases}\label{sec:use-cases}

\begin{table}[htbp]
  \centering
  \caption{Use-case summary}
  \label{tab:use-cases}
  \begin{tabularx}{\textwidth}{@{}l l X l@{}}
    \toprule
    \textbf{ID} & \textbf{Actor} & \textbf{Goal} & \textbf{Requirements} \\
    \midrule
    UC1 & Maya & Install the assistant and reach a first successful answer
      & F1, F2, NF3 \\
    UC2 & Dev & Ask, by voice, whether someone has replied & F3, F4, F13 \\
    UC3 & Dev & Send a reply by voice & F3, F5, F7, F13 \\
    UC4 & Maya & Ask for something new; approve the skill the assistant
      offers to build & F8, F9, F10 \\
    UC5 & Dev & Repeat a previously learned task and see it complete faster
      & F10, NF7 \\
    UC6 & Ines & Inspect an ability's permissions; remove it
      & F10, F11, F18 \\
    UC7 & Ines & Watch the assistant work; review its performance
      measurements & F12, F19 \\
    UC8 & Maya & Reach a login wall and take over herself & F14 \\
    UC9 & Dev & Ask from his phone & F21 \\
    UC10 & Maya & Have an appointment from an email put on her calendar
      & F20, F7, F13 \\
    \bottomrule
  \end{tabularx}
\end{table}

\subsubsection{UC1 --- First run to first answer}\label{sec:uc1}

\begin{itemize}[itemsep=1pt]
  \item \emph{Precondition:} installer obtained; clean machine.
  \item \emph{Trigger:} first launch.
\end{itemize}

\begin{enumerate}[itemsep=1pt]
  \item The assistant checks the machine meets its needs and says what it
    will download and why.
  \item The user accepts the defaults; downloads proceed with visible
    progress and survive interruption.
  \item Everything needed starts without further action; the assistant
    invites a first request.
  \item The user asks a question and receives an answer.
\end{enumerate}

\begin{itemize}[itemsep=1pt]
  \item \emph{Error paths:} machine below requirements $\rightarrow$ the
    missing requirement is named and setup stops cleanly; interrupted
    download $\rightarrow$ resumes on next launch.
\end{itemize}

\subsubsection{UC3 --- Send a reply by voice}\label{sec:uc3}

\begin{itemize}[itemsep=1pt]
  \item \emph{Precondition:} the user has connected their mail account (F14:
    by signing in themselves).
  \item \emph{Trigger:} the user says ``reply to \texttt{<address>} saying
    I'll be there at nine.''
\end{itemize}

\begin{enumerate}[itemsep=1pt]
  \item The transcript is shown; the activity view narrates the steps (F12).
  \item Because the spoken request itself asks for a reply, sending is
    authorised without a second prompt (F7).
  \item The outcome is confirmed from the mailbox itself, and spoken back
    (F13).
\end{enumerate}

\begin{itemize}[itemsep=1pt]
  \item \emph{Variant:} ``draft\ldots'' $\rightarrow$ composed but never
    sent, and the user is told where it is.
  \item \emph{Error path:} the named person is ambiguous in this mailbox
    $\rightarrow$ the assistant declines to guess, lists the candidates, and
    sends nothing.
\end{itemize}

\subsubsection{UC4 --- Approve a new skill}\label{sec:uc4}

\begin{itemize}[itemsep=1pt]
  \item \emph{Trigger:} ``rename all the receipts in this folder by their
    dates.''
\end{itemize}

\begin{enumerate}[itemsep=1pt]
  \item No installed ability covers it; the assistant offers to build one,
    stating what it would do, what it would be allowed to access, and
    roughly how long it will take (F8).
  \item The user approves; progress and test results stream in the activity
    view (F12).
  \item Tests pass $\rightarrow$ the skill is installed, run, and the result
    presented; it now appears among the user's abilities with its origin
    recorded (F9, F10).
  \item Had the user declined, nothing would have been created, and the
    offer lapses on its own (F8).
\end{enumerate}

\begin{itemize}[itemsep=1pt]
  \item \emph{Error path:} attempts exhaust without passing tests
    $\rightarrow$ an honest failure with the reason; nothing is installed
    (F9, F13).
\end{itemize}

\subsection{User flows}\label{sec:user-flows}

\begin{figure}[htbp]
  \centering
  \fbox{\parbox{0.9\textwidth}{\centering\emph{Flow diagram --- exported
    from architecture sources}}}

  \medskip
  \begin{minipage}{0.9\textwidth}
    \small
    \begin{enumerate}[itemsep=1pt]
      \item A user request arrives, typed or spoken, and is matched against
        the assistant's understood abilities.
      \item A question $\rightarrow$ answered from local knowledge; the
        reply is shown or spoken.
      \item An installed ability $\rightarrow$ it is run; the reply is shown
        or spoken.
      \item Mail or calendar $\rightarrow$ the assistant works the account.
        At an irreversible step: if the request authorised it, the assistant
        does it and verifies the outcome; if not, it asks the user ---
        approval leads to the same verified action, decline ends with a
        reply and nothing done.
      \item Nothing covers it $\rightarrow$ the assistant offers to build a
        skill. If approved, it builds and tests: tests pass $\rightarrow$
        install and run, then reply; tests fail $\rightarrow$ an honest
        failure report. If declined or lapsed, the interaction ends with a
        reply and nothing created.
    \end{enumerate}
  \end{minipage}
  \caption{Request flow (every interaction takes this path)}
  \label{fig:request-flow}
\end{figure}

\begin{figure}[htbp]
  \centering
  \fbox{\parbox{0.9\textwidth}{\centering\emph{Flow diagram --- exported
    from architecture sources}}}

  \medskip
  \begin{minipage}{0.9\textwidth}
    \small
    \begin{enumerate}[itemsep=1pt]
      \item First launch.
      \item Machine check --- if the machine is below requirements, the
        missing requirement is named and setup stops cleanly.
      \item Confirm downloads.
      \item Downloads proceed with visible progress and are resumable.
      \item Everything starts.
      \item The user connects mail, signing in themselves.
      \item First request.
    \end{enumerate}
  \end{minipage}
  \caption{First-run flow}
  \label{fig:first-run-flow}
\end{figure}

\begin{figure}[htbp]
  \centering
  \fbox{\parbox{0.9\textwidth}{\centering\emph{Flow diagram --- exported
    from architecture sources}}}

  \medskip
  \begin{minipage}{0.9\textwidth}
    \small
    \begin{enumerate}[itemsep=1pt]
      \item Hold to talk.
      \item On-device transcription.
      \item Transcript shown.
      \item Request flow (Figure~\ref{fig:request-flow}).
      \item On-device speech.
      \item Spoken reply.
    \end{enumerate}
  \end{minipage}
  \caption{Voice flow}
  \label{fig:voice-flow}
\end{figure}

\subsection{Verification mapping (feeds Chapter 5)}\label{sec:verification-mapping}

\begin{table}[htbp]
  \centering
  \caption{Verification mapping}
  \label{tab:verification-mapping}
  \begin{tabularx}{\textwidth}{@{}l X@{}}
    \toprule
    \textbf{Requirement(s)} & \textbf{Verified by} \\
    \midrule
    F4--F7, F13, NF5 & Task suite A of the evaluation protocol, outcomes
      checked against the mailbox itself \\
    F20 & A booking case added to suites A and D, outcome checked against
      the calendar itself \\
    F8, F9 & Consent tests (automated); skill-generation suite C with damage
      rate \\
    F7, F14, NF2 & Safety suite D, including the adversarial page scenario \\
    NF7, F17 & Repeat-cost suite E \\
    F1, NF3 & Scripted clean-machine installation; UAT ($n \geq 4$) with
      SUS \\
    NF1 & Network egress log captured across a full demo script \\
    NF4 & Existing memory-residency measurements \\
    \bottomrule
  \end{tabularx}
\end{table}

\subsection{Decisions log}\label{sec:decisions-log}

\begin{itemize}[itemsep=2pt]
  \item 2026-08-09 --- Phone access (F21) held at COULD.
  \item 2026-08-09 --- General web automation excluded (F25) on measured and
    community evidence, recorded in \S\ref{sec:moscow}'s scope remarks.
  \item 2026-08-09 --- Requirements rewritten solution-agnostically; all
    implementation nouns moved out of this chapter.
  \item 2026-08-10 --- Requirements restated in the department's convention:
    every row ``the application shall / shall not'', prohibitions promoted
    from WON'T to binding MUSTs, non-functional categories removed.
  \item 2026-08-10 --- Calendar booking added (F20, SHOULD): events booked on
    request or from details found in email, behind the F7 approval rules;
    scheduled at the top of Phase~2 with a gated week-4 stretch slot.
\end{itemize}

\end{document}
